\section{Common setup}
\label{sec:setup}

\paragraph{Data and split.} We use CIFAR-10~\cite{krizhevsky2009learning} (50,000 training and 10,000 test images, 10
classes). The 50,000 training images are split \emph{once} into 45,000 train and 5,000 validation images with
scikit-learn's \texttt{train\_test\_split} (stratified by class, \texttt{random\_state}\,=\,\res{splitstate}). The random
state is the integer formed by all digits of the roll number, in order (seed \res{seed}). The split is exactly stratified (4,500/500 per class), shared by all three questions and saved as
\texttt{results/q1\_outputs/split\_indices.npz}. Per-channel normalisation statistics are computed on the 45k training
images only. Q1 and Q2 feed normalised images to the classifiers; the Q3 autoencoders work on $[0,1]$ pixels and
normalise only when passing images to the classifier.

\paragraph{Selection and test protocol.} The validation split is used for every choice: learning rate, early-stopping
epoch, bottleneck, blur strength and probe hyperparameters. Each notebook loads the official test set only in its last
section, after all choices are fixed, evaluates every final model once, and quotes test numbers as final results.
Training and validation curves are reported where a question asks about training behaviour. Every predicted outcome
and decision threshold in this report was written into the notebook before the experiment was run.

\paragraph{Randomness.} Every random choice is derived from the roll-number seed. The split uses the seed directly; all
other randomness draws from NumPy's \texttt{SeedSequence} seeded with the roll-number seed and a descriptive tag (\eg the
pixel permutation, label noise, tile orders, test noise, run seeds and probe seeds). A run's seed index fixes both its
initialisation stream and its mini-batch order, and is shared across models and conditions, which makes comparisons
paired. Unless stated otherwise, results are mean$\pm$sd over seeds.

\paragraph{Training.} All networks are trained with AdamW~\cite{loshchilov2019decoupled} and batch size 128, using
mixed precision (fp16) on the GPU with the whole split held in GPU memory. Classifiers use BatchNorm~\cite{ioffe2015batch}
and ReLU in both families. \textbf{No data augmentation is used anywhere}: crops and flips hard-code exactly the spatial
priors that Q1 and Q2 study. Hardware and software for the reported run: \res{env.device}, Python \res{env.python},
PyTorch \res{env.torch}, torchvision \res{env.torchvision}, CUDA \res{env.cuda}~\cite{paszke2019pytorch}.

\begin{table}[t]
\centering\small\setlength{\tabcolsep}{3pt}
\caption{Classifiers used in Q1 and Q2. MLPs have three hidden layers (Linear--BN--ReLU) of the given width. MLP-M's
width (305) is solved so that its parameter count matches the CNN's within 0.2\%. The CNN is six $3{\times}3$
convolutions (32,64 | 128,128 | 256,256 channels, BN+ReLU) in three stages with $2{\times}2$ max-pooling, followed by
global average pooling and a linear head. MACs: multiply--accumulates per image.}
\label{tab:models}
\begin{tabular}{@{}llrr@{}}
\toprule
Model & Architecture & Parameters & MACs \\
\midrule
MLP-S  & $3072{\to}128^{\times3}{\to}10$  & \res{q1.MLP-S.params}  & \res{q1.MLP-S.macs} \\
MLP-M  & $3072{\to}305^{\times3}{\to}10$  & \res{q1.MLP-M.params}  & \res{q1.MLP-M.macs} \\
MLP-L  & $3072{\to}1024^{\times3}{\to}10$ & \res{q1.MLP-L.params}  & \res{q1.MLP-L.macs} \\
MLP-XL & $3072{\to}2048^{\times3}{\to}10$ & \res{q1.MLP-XL.params} & \res{q1.MLP-XL.macs} \\
CNN    & 6 conv + GAP + linear           & \res{q1.CNN.params}    & \res{q1.CNN.macs} \\
\bottomrule
\end{tabular}
\end{table}
